\documentclass[conference]{IEEEtran}
\usepackage{amsmath,amssymb}
\usepackage{graphicx}
\usepackage{booktabs}
\usepackage{url}

\begin{document}

\title{Do Graph Neural Networks Predict Financial Shock Propagation?\\
A Leakage-Free Evaluation of GATv2 and a Temporal Multi-Market Extension}

% TODO: replace with your name and affiliation
\author{\IEEEauthorblockN{Author Name}
\IEEEauthorblockA{Affiliation \\ email}}

\maketitle

\begin{abstract}
We study whether graph neural networks (GNNs) can predict which stocks make a large move
shortly after another stock experiences an abnormal return. We evaluate two models under a
protocol designed to avoid look-ahead bias: chronological train/validation/test splits with
purge gaps, trailing-only features, dynamic-correlation parameters estimated on training
data only, five random seeds, and non-graph baselines. The first model is a two-layer GATv2
over ten US banks with a one-day horizon. An earlier version of this project reported
0.914 ROC-AUC for it; that figure came from selecting the best of 1,080 configurations on a
12-sample validation set. Under the corrected protocol it reaches ABSTRACT_FIN on held-out
events, which is close to chance. The second model, a multi-market temporal GNN, combines per-snapshot GATv2
layers with transformer attention over twelve weekly graph snapshots of 194 stocks in four
markets, using DCC-GARCH correlations as edge weights. ABSTRACT_ABL We conclude that, for
this label definition, the predictive signal is dominated by each stock's own volatility,
and that the graph and temporal components do not add measurable value.
\end{abstract}

\section{Introduction}
Shocks to one firm can spread to others through common exposures, funding links and
correlated trading. Graph neural networks are a natural fit: stocks are nodes,
relationships are edges, and message passing could model how a shock travels. Financial
prediction is also prone to evaluation leakage: future information enters the features,
events are split randomly in time, or a model is chosen on the same data it is scored on.
These problems can make weak models look strong.

This report makes three contributions. (i) A leakage-free event-based evaluation pipeline
for shock-propagation prediction, with tests that check for look-ahead. (ii) A
re-evaluation of a finance-sector GATv2 whose originally reported performance did not
replicate. (iii) A four-step ablation of a temporal multi-market GNN against
non-graph baselines.

\section{Problem Setup}
Let $P_i(t)$ be the USD close of stock $i$ on trading day $t$ and $r_i(t)$ its daily return.
A \emph{shock} at stock $s$ on day $d$ is a return with $|r_s(d)|/\hat\sigma_s(d-1) > 3$,
where $\hat\sigma_s(d-1)$ is the standard deviation over the trailing 60 days (the multi-market model
additionally requires $|r_s(d)| > 3\%$). At most three shocks are kept per day. For each
event and each other stock $i$, the label is
\begin{equation}
y_i = \mathbb{1}\left[\,\left|P_i(d+h)/P_i(d) - 1\right| > 3\%\,\right].
\end{equation}
Events are split by date into 60\% train, 20\% validation and 20\% test, with an $h$-day gap
between splits. The shock source is excluded from scoring.

\section{Data}
Daily prices come from Yahoo Finance for 2021-10-01 to 2026-09-29. \emph{Finance sector:}
ten large US banks; horizon $h=1$; 88/38/30 train/validation/test events. \emph{Multi-market:}
the 50 largest S\&P~500 stocks, NIFTY~50, and the 50 largest FTSE~100 and Nikkei~225 stocks.
After dropping tickers with missing data, 194 remain. Non-US closes are converted to USD and
aligned on the NYSE calendar. Tokyo, Mumbai and London close before New York on the same
date, so this alignment introduces no look-ahead. The horizon is $h=5$, giving
1162/378/392 events; the test period runs from 2025-11-11 to 2026-09-22.

\section{Models}
\subsection{Graphs and features}
\emph{Finance sector:} documented peer links plus edges from the shock source to every bank.
Each node has 8 features: trailing return statistics, volatility, price and volume ratios, a
shock-source flag and a peer-link flag. \emph{Multi-market:} each node keeps its 10 neighbours
with the largest absolute correlation. Correlations are either Pearson over a trailing
120-day window, or DCC(1,1)-GARCH(1,1)~\cite{engle2002} estimated on training days only and
filtered forward. With 194 assets, the full DCC likelihood drove the news parameter to
$a\approx 0$ (effectively constant correlation). We therefore use the pairwise composite
likelihood~\cite{pakel2021}, which gives $a=0.018$, $b=0.928$. Each node has 14 features,
z-scored with training-period statistics.

\subsection{Architectures}
\emph{StaticGATv2}: two GATv2~\cite{brody2022} layers (2 heads, then 1 head; hidden size 32)
with edge weights in the attention score. \emph{TemporalGATv2}: the same spatial block,
shared across $T=12$ snapshots spaced 5 trading days apart. A pre-LayerNorm transformer
block~\cite{vaswani2017} with sinusoidal positions runs over each node's 12 embeddings; the
last position feeds an MLP head. \emph{Baselines}: GAT, GCN~\cite{kipf2017},
GraphSAGE~\cite{hamilton2017}, a per-node MLP, logistic regression, and a heuristic that
scores each stock by its trailing 20-day volatility.

\subsection{Training}
We use weighted binary cross-entropy (positives weighted by neg/pos) and Adam with gradient
clipping. Early stopping uses validation ROC-AUC. The F1 threshold is chosen on validation
and applied unchanged to test. Every trained model runs with 5 seeds.

\section{Results}
We report test ROC-AUC, PR-AUC (a random scorer achieves the positive base rate) and F1,
as mean $\pm$ standard deviation over seeds.

\begin{table}[t]
\centering
\caption{Finance sector, test set (270 stock-events, base rate 0.093).}
\label{tab:fin}
\begin{tabular}{lccc}
\toprule
Model & ROC-AUC & PR-AUC & F1 \\
\midrule
Volatility heuristic & 0.516 & 0.102 & 0.000 \\
Logistic regression & 0.385 & 0.086 & 0.065 \\
MLP & 0.359 $\pm$ 0.026 & 0.077 $\pm$ 0.007 & 0.053 $\pm$ 0.048 \\
GATv2 & 0.553 $\pm$ 0.022 & 0.114 $\pm$ 0.011 & 0.049 $\pm$ 0.052 \\
GAT & 0.529 $\pm$ 0.023 & 0.108 $\pm$ 0.008 & 0.048 $\pm$ 0.047 \\
GCN & 0.464 $\pm$ 0.032 & 0.089 $\pm$ 0.008 & 0.043 $\pm$ 0.055 \\
GraphSAGE & 0.492 $\pm$ 0.071 & 0.092 $\pm$ 0.014 & 0.021 $\pm$ 0.028 \\
\bottomrule
\end{tabular}
\end{table}

\begin{table}[t]
\centering
\caption{Multi-market ablation, test set (S\&P 500 targets, base rate ABL_BASE).}
\label{tab:abl}
\begin{tabular}{lccc}
\toprule
Model & ROC-AUC & PR-AUC & F1 \\
\midrule
ABLATION_ROWS
\bottomrule
\end{tabular}
\end{table}

\textbf{Finance sector} (Table~\ref{tab:fin}). GATv2 ranks first, but every model is close
to chance. With only 25 positive labels in the test set, these differences are not
reliable. The non-graph models score below 0.5, which suggests that the relationships
learned from 2022--2024 reversed in the test period.

\textbf{Multi-market ablation} (Table~\ref{tab:abl}). ABLATION_DISCUSSION

\section{Lessons from the Original Evaluation}
The originally reported 0.914 ROC-AUC had four causes. (1) It was the maximum over 1,080
configurations, each scored on 12 validation samples. (2) Events were ordered by ticker
before a ``temporal'' split, so training and validation overlapped in time. (3) Shock
z-scores used full-sample statistics. (4) The reported number was also the model-selection
number. Retraining the same configuration gave 0.69 and 0.76 at the time. Under the protocol
above it gives 0.553.

\section{Limitations}
The test sets are small and come from a single market regime. The label (absolute move
above 3\%) is largely explained by each stock's own volatility; a label based on abnormal
returns would isolate propagation better. The universe consists of current large caps, so
it has survivorship bias. Multi-market model hyperparameters were not tuned. Attention weights are
diagnostics, not causal attributions.

\section{Conclusion}
Under a leakage-free protocol, neither a finance-sector GATv2 nor a temporal multi-market
GNN outperforms a trailing-volatility heuristic at predicting which stocks move after a
shock. The evaluation pipeline, including its look-ahead tests, is the more durable
contribution.

\begin{thebibliography}{9}
\bibitem{brody2022} S.~Brody, U.~Alon, and E.~Yahav, ``How attentive are graph attention networks?'' in \emph{ICLR}, 2022.
\bibitem{engle2002} R.~Engle, ``Dynamic conditional correlation: A simple class of multivariate generalized autoregressive conditional heteroskedasticity models,'' \emph{J. Business \& Economic Statistics}, vol.~20, no.~3, pp.~339--350, 2002.
\bibitem{pakel2021} C.~Pakel, N.~Shephard, K.~Sheppard, and R.~Engle, ``Fitting vast dimensional time-varying covariance models,'' \emph{J. Business \& Economic Statistics}, vol.~39, no.~3, pp.~652--668, 2021.
\bibitem{vaswani2017} A.~Vaswani \emph{et al.}, ``Attention is all you need,'' in \emph{NeurIPS}, 2017.
\bibitem{kipf2017} T.~Kipf and M.~Welling, ``Semi-supervised classification with graph convolutional networks,'' in \emph{ICLR}, 2017.
\bibitem{hamilton2017} W.~Hamilton, R.~Ying, and J.~Leskovec, ``Inductive representation learning on large graphs,'' in \emph{NeurIPS}, 2017.
\end{thebibliography}

\end{document}
